\honest{Argument Screens Off Authority}{Eliezer Yudkowsky}{2007}

A famous geologist and a fourteen-year-old juvenile delinquent each make an odd claim about rocks. Arthur believes the geologist at 90\% and the boy at 10\%. That is authority. Two men make odd claims about physics, one with a careful argument and one with leaps of faith, and Arthur believes them at 90\% and 10\%. That is argument.

Now swap. The geologist and the boy give equally good technical cases, and the two physicists turn out to have equal credentials. The geologist's advantage mostly vanishes, and the careful arguer's advantage stays. So once you have the arguments, authority adds little. This holds only if Arthur can understand all the technical arguments; ``otherwise they're just impressive noises.'' Otherwise, I grant, belief ``depends a great deal on his authority.'' \nb{People rely on authority mostly when they cannot follow the argument, so the principle covers the cases where they rely on it least.}

I also require that each speaker give ``the best argument he could give,'' including every step and every piece of support he used. \nb{Outside mathematics, almost no real argument meets this.}

A novice would conclude that authority and argument are different kinds of evidence, since both give 90\% against 10\% yet combine differently. I show that probability theory handles this. If one fact makes a hypothesis 90\% likely and another makes it 9\% likely, the probability given both cannot be computed from those two numbers, because the facts may depend on each other. Take a sidewalk that is slippery while a sprinkler runs, and a sprinkler that runs 10\% of the night. Given that the sprinkler is on, the sidewalk is slippery with probability 90\%; given that it is night, 9\%; given both, 90\% again. Once you know whether the sprinkler is on, knowing whether it is night tells you nothing more about the sidewalk. This is called screening off. In a causal graph the arrows run from night to sprinkler to sidewalk, and their direction matters. Two dice that add to a sum are the opposite case: the first die tells you nothing about the second until you learn the sum, and then it tells you everything. The rule for reading such independences off a graph is called D-separation, and the books to read are Judea Pearl's. \nb{The example is built so that night affects the sidewalk only through the sprinkler: there is no rain and no dew.}

Then my diagram for experts. The truth causes good arguments, and good arguments cause experts to believe. So once you know the argument, the expert's belief adds nothing about the truth. The diagram has no other path from the truth to the expert. None of this needs a second kind of probability; it is ordinary probability, drawn so the independences are easy to see. I draw it that way; I do not argue for it. Beside it I write ``(In theory!)''

In my closing paragraphs I supply the other paths myself. Good experts know counterevidence they did not mention. Judging a step may need ``intuitions you can't duplicate without the same thirty years of experience.'' \nb{Each of these means the argument the listener hears is not all of what the expert knows. Hearing it then does not fully screen off the expert. The post calls this a matter of practice, and keeps its title.}

I grant that what E. T. Jaynes says about probability deserves slightly more credence than the same words from me. ``But this slight strength of authority is only ceteris paribus.'' I once found an error in one of his books, ``because algebra trumps authority.'' \nb{That is true, for algebra.}
